\section{Discussão e Conclusão}

A realização experimental e posterior análise foram elucidativas, na medida que permitiram ao grupo de trabalho perceber como ocorre a carga e a descarga de um condensador. Nestes processos, os tempos de carga e descarga dependem das resistências do circuito, de como elas se encontram associadas e do condensador, conforme a \cref{eq:tauc}.
Os valores obtidos para as constantes do tempo $\tau_c = 10,30\,$s e $\tau_d = 12,25$\,s revelaram-se superiores aos valores teóricos $\tau_c = 6\,s$ e $\tau_d = 10\,s$.
A este desvio pode-se atribuir pelo menos um fator enraizado em erros experimentais, como a utilização do valor ideal das resistências em vez do seu valor real, este último tangível através de um ohmímetro. É também importante referir que este fator pode ser agravado pelas tolerâncias das resistências e do condensador, que podem aumentar o tempo de carga e descarga.






